\section{Conclusion and Discussion}

We presented a reproducibility-focused study of calibration and frozen-feature
adaptation for AutoShot. The selected configuration trains a BCE one-hot head
on cached features and applies temperature scaling and Gaussian smoothing
chosen by source-stratified, validation-only cross-validation. It reaches F1
scores of \PaperASVTwoShotFOne{} on SHOT, \PaperASVTwoBBCFOne{} on BBC, and
\PaperASVTwoClipFOne{} on ClipShots for the reported seed.

The ablation clarifies the source of the gain: Focal Loss and many-hot
supervision do not improve the BCE control, whereas Gaussian smoothing changes
boundary decisions and temperature scaling improves probability diagnostics.
Paired bootstrap intervals support gains on SHOT and ClipShots but not BBC.
Class-balanced calibration error also shows why very small aggregate ECE values
should not be interpreted as uniformly calibrated boundary probabilities.

The revised implementation records exact cache identities, split manifests,
frozen selection decisions, and artifact checksums. Completing the scheduled
three-seed study requires restoring the original AutoShot base checkpoint;
until then, optimization variability remains unresolved. Further work should
also reserve an independent calibration split and address cross-domain
precision on ClipShots.
